\section{Assembler}
\label{sec:toolchain-assembler}

The \texttt{viv32-asm} tool translates a VIV-32 assembly source into a relocatable object
file, suitable for linking with \texttt{viv32-ld}.
Instruction encodings and architectural instruction semantics are defined by the VIV-32
architecture specification; this section only describes assembler-level syntax, labels,
sections, directives, and symbol resolution.

\subsection{Source format}

Assembly source is line-oriented. Each non-empty line may contain a section declaration,
a label, a directive, an instruction, or a comment.
A label may appear on its own line or before an instruction or directive on the same line.

Comments begin with \texttt{;} and continue to the end of the line.

\begin{verbatim}
; This is a comment
add $1, $2, $3
\end{verbatim}

Whitespace is insignificant except where required to separate tokens.

\subsection{Case sensitivity}

Assembly source is case-insensitive. Instruction mnemonics, directives, register names,
control register names, section names, and symbol names may be written in upper case,
lower case, or mixed case.

The assembler normalises symbol names before recording definitions and references.
Consequently, the following labels refer to the same symbol name:

\begin{verbatim}
loop:
LOOP:
Loop:
\end{verbatim}

Defining more than one label whose names differ only by case is an error.

\subsection{Registers}

General-purpose registers are written with a leading dollar sign:

\begin{verbatim}
$0
$1
$15
\end{verbatim}

Control registers are written with a leading percent sign and their architectural name:

\begin{verbatim}
%pc
%sr
%evbase
%ecause
%epc
%esr
%edata
\end{verbatim}

\subsection{Sections}

The supported sections are:

\begin{description}
    \item[\texttt{.text}] Executable instructions.
    \item[\texttt{.rodata}] Emitted read-only data, such as strings and lookup tables.
    Labels in this section define address-bearing symbols.
    \item[\texttt{.data}] Emitted initialised mutable data.
    Labels in this section define address-bearing symbols.
    \item[\texttt{.bss}] Zero-initialised reserved storage.
    Labels in this section define address-bearing symbols.
\end{description}

Constants declared with \texttt{.const} are not part of any section and are not emitted.

In the initial relocatable object format, sections primarily classify emitted content.
The linker preserves object order as provided on the command line. Future linkers may
use section information for more sophisticated layout.

\subsection{Labels, symbols, and name mangling}

A label defines a symbol at the current assembly address. Labels are written as an
identifier followed by a colon:

\begin{verbatim}
start:
loop:
message:
\end{verbatim}

Labels may be used as operands where an address or relative branch target is required.
The assembler marks label declarations, use, and mangling. Labels are mangled by
default, using a file-unique identifier.

Shared labels can be indicated with \texttt{.nomangle} to prevent mangling, allowing
multiple files access to the labeled item.

Symbol names must begin with a letter or underscore, followed by zero or more letters,
digits, or underscores.

\begin{verbatim}
valid_label:
_valid_label:
label_1:

\end{verbatim}

\subsubsection{\texttt{.nomangle}}

By default, labels are local to the assembled object. This allows separate assembly
files to use common internal label names such as \texttt{loop}, \texttt{done}, or
\texttt{error} without creating linker-level name conflicts.

The \texttt{.nomangle} directive marks a symbol name as externally visible. A
non-mangled symbol is placed in the global linker namespace and may be referenced by
other assembled objects.

The \texttt{.nomangle} directive must appear before the first use or definition of the
named symbol in the source file. This allows the assembler to determine whether each
symbol reference is object-local or linker-global when it is first encountered.

\begin{verbatim}
.nomangle print_string

print_string:
    ; function body
\end{verbatim}

A file that references an externally defined symbol should also mark that symbol as
non-mangled:

\begin{verbatim}
.nomangle print_string

    call print_string
\end{verbatim}

The directive does not emit bytes and does not define storage. It only controls how the
assembler records the symbol name in the relocatable object. Whether the symbol is
defined in the current object or resolved from another object is determined by the
linker.

\texttt{.nomangle} applies to address-bearing symbols only. Constants declared with
\texttt{.const} are assembler-local substitutions and are not exported to the linker.

\subsection{Constants and address-bearing symbols}

The assembler distinguishes constants from address-bearing symbols.

Constants are assembler-time numeric substitutions. They do not occupy space in the
output image and do not define storage. A constant is declared with \texttt{.const}:

\begin{verbatim}
.const SERIAL_DATA,    0xFFFF0000
.const SERIAL_CONTROL, 0xFFFF0004
.const SERIAL_ENABLE,  0x00000001
.const SERIAL_TX_EN,   0x00000002
\end{verbatim}

Constant values must be unsigned hexadecimal integer literals. When a constant is used
as an operand, the assembler substitutes its value directly into the instruction or
directive being encoded.

\begin{verbatim}
lli $2, SERIAL_ENABLE
sw  $2, [$1, 0]
\end{verbatim}

A label, by contrast, defines an address-bearing symbol. Labels mark locations in
sections such as \texttt{.text}, \texttt{.rodata}, \texttt{.data}, or \texttt{.bss}.
When a label is used, the assembler or linker resolves it as an address or as a
PC-relative offset, depending on the instruction or directive.

\begin{verbatim}
.rodata:
message:
    .asciz "Hello World!\n"

.text:
    li $2, message
\end{verbatim}

In this example, \texttt{message} is not a constant. It refers to the address of data
stored in the output image.

\subsection{Directives}

Assembler directives begin with a period.

\subsubsection{\texttt{.org}}

The \texttt{.org} directive sets the current object-local assembly offset.

\begin{verbatim}
.org 0x00000100
\end{verbatim}

If the new offset is greater than the current object-local offset, the assembler emits
zero padding within the relocatable object. If the new offset is less than the current
object-local offset, assembly fails. Final executable addresses are assigned by
\texttt{viv32-ld}.

\subsubsection{\texttt{.byte}}

The \texttt{.byte} directive emits one or more byte values.

\begin{verbatim}
.byte 0x48, 0x65, 0x6C, 0x6C, 0x6F
\end{verbatim}

Each value must fit in 8 bits.

\subsubsection{\texttt{.word}}

The \texttt{.word} directive emits one or more 32-bit words in big-endian byte order.

\begin{verbatim}
.word 0x12345678
.word start
\end{verbatim}

Each value must fit in 32 bits after symbol resolution.

\subsubsection{\texttt{.ascii}}

The \texttt{.ascii} directive emits the bytes of a string literal without a terminating
zero byte.

\begin{verbatim}
.ascii "Hello"
\end{verbatim}

\subsubsection{\texttt{.asciz}}

The \texttt{.asciz} directive emits the bytes of a string literal followed by a zero
terminator.

\begin{verbatim}
.asciz "Hello World!\n"
\end{verbatim}

\subsubsection{\texttt{.space}}

The \texttt{.space} directive reserves a number of zero-filled bytes.

\begin{verbatim}
.space 64
\end{verbatim}

\subsection{Instruction operands}

The canonical instruction format uses comma-separated operands:

\begin{verbatim}
add  $1, $2, $3
addi $1, $2, 4
\end{verbatim}

Memory operands are written using square brackets. The base register is listed first,
followed by the signed byte offset.

\begin{verbatim}
lw $1, [$2, 0]
sw $3, [$4, -4]
\end{verbatim}

The brackets are a source-level aid for readability. The underlying instruction format
contains the register and offset fields directly.

\subsection{Immediates and numeric literals}

Decimal and hexadecimal integer literals are supported.

\begin{verbatim}
42
-4
0x1234
0xFFFF0000
\end{verbatim}

Hexadecimal literals use the \texttt{0x} prefix. Signed offsets are normally written in
decimal source form, while bit patterns, masks, addresses, and immediate-construction
values are normally written in hexadecimal.

\subsection{Branch and jump symbols}

Branch and jump instructions may use either numeric offsets or labels.

\begin{verbatim}
loop:
    addi $1, $1, -1
    b.ne $1, $0, loop
\end{verbatim}

VIV-32 PC-relative control-flow offsets are relative to the instruction following the
branch or jump.
When a label is used, the assembler records the reference. If the target address is not
final at assembly time, the object contains a relocation record. The linker computes the
final PC-relative offset.

The resulting offset must be aligned and must fit in the target instruction's encoded
offset field. If it does not, assembly fails.

\subsection{Pseudo-instructions}

Pseudo-instructions are expanded before object-local offsets are assigned, so each
emitted item has a fixed size before relocation records are produced.

Initial pseudo-instructions may include:

\begin{description}
    \item[\texttt{li}] Load a 32-bit immediate value or symbol address.
    \item[\texttt{la}] Load the address of a symbol.
    \item[\texttt{ret}] Return using the ABI-defined return-address register.
\end{description}

Pseudo-instruction expansion has fixed size. If a pseudo-instruction cannot be expanded
unambiguously, assembly fails.

\subsection{Example}

\begin{verbatim}
.text:
.org 0x0                       ; The reset interrupt, also program entry.
    jmp     _start             ; Jump past the interrupt vector table.

.org 0x7C                      ; Final instruction of interrupt vector table,
    halt                       ; All interrupts effectively nop until halt.
_start:
    lli     $1, 0x0003         ; Serial control mask to enable TX.
    li      $2, 0xFFFF0200     ; Serial MMIO base address.
    sw      $1, [$2, 4]        ; Apply serial control mask.

loop:
    lbu     $1, [$3, message]  ; Read the first byte of message, using $3==0
                               ; as first index. Note the reversal of base and
                               ; offset in this example enables fewer
                               ; instructions.
    b.eq    $1, $0, _end       ; Exit loop at end of message.
    sb      $1, [$2, 0]        ; Send the byte to serial data buffer.
    addi    $3, $3, 1          ; Increment our index $3.
    jmp     loop

_end:
    halt

.rodata:
.org 0x0100
message: .asciz "Hello World!\n"
\end{verbatim}
